% !TEX TS-program = lualatex
% !TEX encoding = UTF-8 Unicode
\documentclass[ngerman,11pt]{article}
\usepackage{fontspec}
\usepackage[ngerman]{babel}
\usepackage[ngerman]{datetime2}
\usepackage[a4paper, total={16cm, 25cm}]{geometry}
\usepackage{amsmath}
\usepackage{amsfonts}
\usepackage{amssymb}
\usepackage{multicol}
\usepackage{tcolorbox}
\usepackage{cancel}
\usepackage{enumitem}
\usepackage{url}
\usepackage[colorlinks=true, linkcolor=blue, urlcolor=blue]{hyperref}
\usepackage{listings}
\usepackage{array}

\lstset{
    language=Python,
    basicstyle=\ttfamily\small,
    keywordstyle=\color{blue},
    stringstyle=\color{red},
    commentstyle=\color{gray},
    breaklines=true,
    frame=single,
    framesep=1mm,
    framextopmargin=1mm,
    framexbottommargin=1mm,
    showspaces=false,
    showstringspaces=false
}

\lstdefinestyle{inline}{
    basicstyle=\ttfamily\normalsize,
    breaklines=false,
    frame=none
}
\date{\DTMgermanmonthname{\month} \the\year}

\begin{document}
\hrule
\begin{center}
{\large\bf Python: Übungsaufgaben zu Datentypen}\\[-0mm]
Studienkolleg der TU Berlin\hfill Till Zoppke\\
\makeatletter Informatik \hfill \@date\\[2mm] \makeatother
\end{center}
\hrule

% Aufgabe 1
\section*{Aufgabe 1: Fehlerhafte Ausdrücke}
Erklären Sie für die folgenden Ausdrücke jeweils, warum Python sie nicht auswerten kann und einen Fehler anzeigt.

\lstinline[style=inline]{>>> [1, 2] + (3, 4)}

\vspace{0.7cm}

\lstinline[style=inline]{>>> 7, 0 % 0, 7}

\vspace{0.7cm}

\lstinline[style=inline]{>>> "abcdefg"[7]}

\vspace{0.7cm}

\lstinline[style=inline]{>>> sum(1,2,3)}

\vspace{0.7cm}

\lstinline[style=inline]{>>> "'"'}

\vspace{0.7cm}

\lstinline[style=inline]{>>> 7 % len([])}

\vspace{0.7cm}

% Aufgabe 2
\section*{Aufgabe 2: Datentypen bestimmen}
Die folgenden Ausdrücke können von Python ausgewertet werden. Kreuzen Sie jeweils an, welchen Datentyp das Ergebnis der Auswertung hat.

\begin{center}
\begin{tabular}{|l|>{\centering}p{1.1cm}|>{\centering}p{1.1cm}|>{\centering}p{1.1cm}|>{\centering}p{1.1cm}|>{\centering}p{1.1cm}|>{\centering\arraybackslash}p{1.1cm}|}
\hline
\textbf{Ausdruck} & \textbf{int} & \textbf{float} & \textbf{string} & \textbf{bool} & \textbf{list} & \textbf{tuple} \\
\hline
\lstinline[style=inline]|7.0 // 2.0| & & & & & & \\[0.3cm]
\hline
\lstinline[style=inline]|[7] * 0| & & & & & & \\[0.3cm]
\hline
\lstinline[style=inline]|"a > b"| & & & & & & \\[0.3cm]
\hline
\lstinline[style=inline]|(1,2)[0:1]| & & & & & & \\[0.3cm]
\hline
\lstinline[style=inline]|ord('A')| & & & & & & \\[0.3cm]
\hline
\lstinline[style=inline]|sum([not str(7)])| & & & & & & \\[0.3cm]
\hline
\lstinline[style=inline]|round(3.5) - 1| & & & & & & \\[0.3cm]
\hline
\lstinline[style=inline]|all(':)')| & & & & & & \\[0.3cm]
\hline
\lstinline[style=inline]|4 ** 4| & & & & & & \\[0.3cm]
\hline
\lstinline[style=inline]|str([1,2,3])[0]| & & & & & & \\[0.3cm]
\hline
\lstinline[style=inline]|[]+[False]| & & & & & & \\[0.3cm]
\hline
\lstinline[style=inline]|(1,1)*2| & & & & & & \\[0.3cm]
\hline
\end{tabular}
\end{center}

% Aufgabe 3
\section*{Aufgabe 3: Ausdrücke auswerten}
Wie werden die folgenden Python-Ausdrücke in der interaktiven Konsole ausgewertet? Schreiben Sie hinter jeden Ausdruck das Ergebnis, das die Python Konsole ausgeben würde.

\begin{lstlisting}

>>> "'"

>>> [3,2][1]

>>> all([7 > 1, 7 % 1])

>>> chr(ord('7'))

>>> 'A' or 'B'

>>> 22 / 22

>>> (max(1, min(2, 3), 4), 5)

>>> int(bool(list([])))

>>> 7 // 3 // -1

>>> [] or [0] and [0,0]

>>> (2,2) != (2) + (2)

>>> round (0.123 ** 1, 1)

>>> [len(tuple('123'))]

\end{lstlisting}

\end{document}
